\documentclass[11pt,a4paper,fontset=none]{ctexart}
\usepackage[margin=2.1cm]{geometry}
\usepackage{hyperref}
\usepackage{xurl}
\usepackage{booktabs}
\usepackage{array}
\usepackage{longtable}
\usepackage{enumitem}
\usepackage{tikz}
\usetikzlibrary{arrows.meta,positioning,shapes.geometric,fit}

\setlength{\parindent}{2em}
\setlength{\parskip}{0.25em}
\setlength{\emergencystretch}{3em}
\setlist{itemsep=0.25em,topsep=0.35em}

\setCJKmainfont[
  ItalicFont={Songti SC},
  BoldItalicFont={Songti SC}
]{Songti SC}
\setCJKsansfont{Hiragino Sans GB}
\setCJKmonofont{Heiti SC}

\hypersetup{
  colorlinks=true,
  linkcolor=black,
  urlcolor=blue,
  pdftitle={班主任.Skill：面向班级治理的多 Skill 工作台},
  pdfauthor={YZDame}
}

\title{\textbf{班主任.Skill：面向班级治理的多 Skill 工作台}}
\author{YZDame\thanks{\href{https://github.com/YZDame}{https://github.com/YZDame}}}
\date{\today}

\begin{document}

\maketitle

\begin{abstract}
班主任工作同时包含学生主档、成绩和表现记录、家校沟通、班务安排以及面向家长和学校的文件产物。若只把这些工作写成一段角色提示词，数据难以复用，写入也缺少可检查的边界。本文介绍开源项目\textbf{班主任.Skill}（\texttt{headteacher-skill}）的 v3 设计：将工作拆成 5 个业务 Skills 和 3 个后端适配 Skills，以统一数据协议连接飞书多维表格、Notion 与 Obsidian，并把 Word、Excel、PPT 视为可追溯的下游产物。系统要求在写入、迁移和删除前提供预览并获得确认。当前飞书是首个优先适配器，Notion 与 Obsidian 提供映射和 fixture；真实远程写入仍取决于用户配置的连接器和凭据。本文是工程设计说明，不宣称已经完成课堂效果评估。
\end{abstract}

\section{问题范围与设计目标}

班主任日常工作有三个相互关联的特征。第一，学生主档相对稳定，成绩、课堂表现、纪律和家校沟通却会持续增加。第二，座位、值日和班委属于带有生效时间的安排，不能简单覆盖旧记录。第三，教师最终需要的是可以直接使用的通知、表格和演示文稿，而不是一堆只停留在聊天窗口里的答案。

因此，本项目把目标限定为四件可检查的事情：

\begin{enumerate}
  \item 让 Agent 能发现并调用职责清晰的 Skills，而不是依赖一段不可拆分的长提示词；
  \item 用后端无关的数据协议保存稳定 ID、时间、来源、可见范围和敏感级别；
  \item 让飞书多维表格、Notion 和 Obsidian 通过适配器承载同一套概念模型；
  \item 在生成或写入前给出预览，写入后保留版本、来源和产物登记信息。
\end{enumerate}

仓库提供的是可安装的工程组件和校验脚本，不包含学生个人数据、账号 Token 或远程服务的默认连接。项目文档中的“支持”指已有协议、映射或 fixture；是否能够实时读写，必须以当前 Agent 的连接器、权限和本地工具检查结果为准。

\section{总体架构}

\subsection{从单一 Skill 到插件包}

v3 使用一个插件清单承载多个可独立发现的 \texttt{SKILL.md}。业务 Skills 处理班级任务，适配 Skills 处理后端差异，\texttt{headteacher-workbench} 负责初始化和路由。这样做可以让用户只安装需要的功能，也可以在 Codex 等支持插件的 Agent 中一次安装完整包；Claude Code、DSH、豆包工作台等环境则可以按各自约定导入同一组完整 Skill 文件夹。

\subsection{请求、协议与后端}

一条请求先由 Agent 识别任务，再进入业务 Skill。业务 Skill 只依赖统一数据协议，不直接假设某个后端的字段名称；适配器负责把协议映射到具体服务。文件生成位于数据链路下游，产物必须记录其来源。图~\ref{fig:architecture} 用普通用户可以理解的方式展示这条链路。

\begin{figure}[htbp]
  \centering
  \begin{tikzpicture}[
      node distance=5mm and 4mm,
      box/.style={rounded corners=2pt, draw=black!65, thick, align=center, minimum height=13mm, text width=2.25cm, font=\small},
      wide/.style={box, text width=2.75cm},
      contract/.style={box, text width=2.45cm, fill=yellow!20, draw=orange!65!black},
      arrow/.style={-{Latex[length=2.5mm]}, thick}
    ]
    \node[box, fill=blue!10, draw=blue!55!black] (request) {教师\\自然语言请求};
    \node[box, fill=blue!10, draw=blue!55!black, right=of request] (agent) {Agent\\识别与路由};
    \node[wide, fill=green!12, draw=green!50!black, right=of agent] (skills) {业务 Skills\\\texttt{class-data} · 成长\\班务 · 产物};
    \node[contract, right=of skills] (contract) {统一数据协议\\Schema + 安全规则};
    \node[wide, fill=orange!14, draw=orange!65!black, right=of contract] (outputs) {适配器与输出\\飞书 · Notion\\· Obsidian\\Office 产物};
    \draw[arrow] (request) -- (agent);
    \draw[arrow] (agent) -- (skills);
    \draw[arrow] (skills) -- (contract);
    \draw[arrow] (contract) -- (outputs);
  \end{tikzpicture}
  \caption{班主任.Skill 的运行链路：请求先经过路由和业务处理，再通过统一协议连接后端或生成下游产物。}
  \label{fig:architecture}
\end{figure}

\section{Skill 组成与职责}

\subsection{业务 Skills}

五个业务 Skills 形成从初始化到产物生成的最小工作链路。每个目录都包含可独立发现的 \texttt{SKILL.md}，并在需要时引用仓库根目录的协议和脚本。

\begin{table}[htbp]
  \centering
  \small
  \caption{业务 Skills 及其职责。}
  \begin{tabular}{p{3.3cm} p{4.2cm} p{7.0cm}}
    \toprule
    Skill & 主要输入 & 主要输出或动作 \\
    \midrule
    \shortstack[l]{\texttt{headteacher-}\\\texttt{workbench}} & 后端选择、初始化请求、环境状态 & 检查环境，选择单一事实来源，并把任务路由到其他 Skills \\
    \texttt{class-data} & 名单、考试、成绩、事件和迁移材料 & 规范化记录，执行导入、追加、查询、更新或 tombstone 预览 \\
    \texttt{student-growth} & 观察、谈话、干预、家校沟通 & 可追溯的成长事件和纵向时间线，区分事实、解释和下一步 \\
    \texttt{class-operations} & 学生属性、座位或值日约束、有效日期 & 带冲突说明的安排预览，以及版本化的 assignment 记录 \\
    \texttt{class-artifacts} & 规范记录、受众、日期范围和模板规则 & Word、Excel、PPT 等下游产物，并登记其来源记录 \\
    \bottomrule
  \end{tabular}
\end{table}

\subsection{后端适配 Skills}

\texttt{feishu-adapter}、\texttt{notion-adapter} 和 \texttt{obsidian-adapter} 共享同一组概念操作：检查前提条件、认证、初始化 Schema、读取实体、写入变更集、生成视图、登记产物和说明限制。业务 Skill 不应绕过这些入口直接拼接后端字段。

\section{统一数据协议}

\subsection{协议外壳与记录不变量}

协议版本为 1.0。每次请求或变更集都带有下表字段；其中 \texttt{dry\_run} 用于明确只预览不写入，\texttt{request\_id} 用于追踪一次操作。

\begin{table}[htbp]
  \centering
  \small
  \caption{协议外壳的关键字段。}
  \begin{tabular}{p{3.2cm} p{10.8cm}}
    \toprule
    字段 & 含义 \\
    \midrule
    \texttt{protocol\_version} & 当前固定为 \texttt{1.0}，用于判断协议兼容性 \\
    \texttt{operation} & \texttt{import}、\texttt{append}、\texttt{upsert}、\texttt{query}、\texttt{tombstone} 或 \texttt{export} \\
    \texttt{request\_id} / \texttt{workspace\_id} & 分别标识请求和班级工作台 \\
    \texttt{actor} / \texttt{emitted\_at} & 记录操作者来源和带时区的发出时间 \\
    \texttt{dry\_run} / \texttt{records} & 标记是否只预览，并携带规范化后的记录数组 \\
    \bottomrule
  \end{tabular}
\end{table}

每条记录至少包含 \texttt{entity\_type}\allowbreak、\texttt{entity\_id}\allowbreak、\texttt{class\_id}\allowbreak、\texttt{revision}\allowbreak、\texttt{source}\allowbreak、\texttt{visibility}\allowbreak、\texttt{sensitivity} 和 \texttt{payload}。实体类型覆盖 \texttt{class\_profile}\allowbreak、\texttt{student\_master}\allowbreak、\texttt{exam\_batch}\allowbreak、\texttt{score\_detail}\allowbreak、\texttt{growth\_event}\allowbreak、\texttt{parent\_contact}\allowbreak、\texttt{seat\_assignment}\allowbreak、\texttt{duty\_assignment}\allowbreak、\texttt{committee\_assignment} 和 \texttt{artifact\_record}。稳定 ID 只在同一工作台内保证不变；时间使用带明确偏移量的 ISO 8601；无法识别的字段保留在 \texttt{payload} 或报告为警告，不能静默丢弃。

\subsection{写入控制与可追溯性}

重复提交同一 \texttt{entity\_id} 和 \texttt{revision} 应保持幂等。删除通过 tombstone 表示，不能把学生记录直接硬删除。迁移前先检查已有后端，新增字段尽量采用可回滚的加法变更。电话、地址和证件号码默认受到限制；分享给家长或学校的产物只能带任务所需字段。

图~\ref{fig:writeflow} 展示一次写入或迁移请求的最短路径。用户可以在确认节点选择停止，或者保持 \texttt{dry\_run} 状态继续查看预览。

\begin{figure}[htbp]
  \centering
  \begin{tikzpicture}[
      node distance=6mm and 7mm,
      flow/.style={rounded corners=2pt, draw=black!65, thick, align=center, minimum height=11mm, text width=2.35cm, font=\small},
      decision/.style={diamond, draw=black!65, thick, aspect=2.1, align=center, inner sep=1pt, text width=1.8cm, font=\small},
      arrow/.style={-{Latex[length=2.5mm]}, thick}
    ]
    \node[flow, fill=blue!10] (input) {提出任务};
    \node[flow, fill=green!12, right=of input] (normalize) {识别实体\\规范化字段};
    \node[flow, fill=yellow!20, right=of normalize] (preview) {生成变更\\预览};
    \node[decision, fill=orange!18, right=of preview] (confirm) {确认写入？};
    \node[flow, fill=orange!14, right=of confirm] (execute) {调用适配器\\查询或写入};
    \node[flow, fill=gray!12, below=of confirm] (stop) {补充信息\\或结束 dry-run};
    \node[flow, fill=purple!12, below=of execute] (record) {返回记录\\登记产物};
    \draw[arrow] (input) -- (normalize);
    \draw[arrow] (normalize) -- (preview);
    \draw[arrow] (preview) -- (confirm);
    \draw[arrow] (confirm) -- node[above,font=\scriptsize]{是} (execute);
    \draw[arrow] (confirm) -- node[right,font=\scriptsize]{否} (stop);
    \draw[arrow] (execute) -- (record);
  \end{tikzpicture}
  \caption{写入、迁移和删除都经过预览与确认；查询可以在预览后直接返回，不必改变数据。}
  \label{fig:writeflow}
\end{figure}

\section{后端适配与当前状态}

工作台选择一个后端作为事实来源。适配器负责投影和读取，不自动在多个后端之间同步同一条记录。这样可以避免同一学生在多个系统中出现无法判断的冲突，也让迁移成为一项显式、可审计的操作。

\begin{table}[htbp]
  \centering
  \small
  \caption{当前后端能力边界。}
  \begin{tabular}{p{3.0cm} p{4.0cm} p{7.2cm}}
    \toprule
    后端 & 仓库内的实现 & 仍需的运行条件 \\
    \midrule
    飞书多维表格 & 首个优先适配器，包含连接路由、Schema 映射和迁移检查 & 用户提供飞书凭据，并配置可用的 API/MCP 或 \texttt{lark-cli}；真实写入须经过预览和确认 \\
    Notion & 适配器骨架、字段映射和 fixture & 已连接的 Notion API/MCP；仓库不宣称已经达到飞书的实时能力 \\
    Obsidian & Markdown、YAML frontmatter 和 Bases 的本地投影适配器 & 指定本地 vault，可选 Obsidian CLI；它保持本地优先，不等同于远程数据库 \\
    \bottomrule
  \end{tabular}
\end{table}

这种状态描述把“协议已经定义”“映射已经提供”“连接器已经可用”分开。Agent 在运行时应先执行环境检查，再向用户说明缺少的工具或权限，而不能根据仓库文件推断已经拥有远程访问能力。

\section{典型使用方式}

用户以自然语言提出任务，Agent 先确定实体、时间范围、受众和目标后端，再调用相应的业务 Skill。下表给出不依赖具体班级数据的示例。

\begin{table}[htbp]
  \centering
  \small
  \caption{自然语言任务与执行链路。}
  \begin{tabular}{p{5.4cm} p{4.0cm} p{4.8cm}}
    \toprule
    用户请求示例 & 主要 Skills & 预期动作 \\
    \midrule
    “把名单和历史成绩导入工作台。” & \texttt{class-data} + 后端适配器 & 规范化后生成 upsert 预览，确认后写入 \\
    “记录今天的课堂表现和一次家校沟通。” & \texttt{student-growth} + \texttt{class-data} & 分别形成 growth event 与 parent contact，保留日期和来源 \\
    “按身高和分组约束重排座位。” & \texttt{class-operations} + \texttt{class-data} & 展示硬约束、冲突和未安排学生，确认后生成带有效日期的安排 \\
    “根据最近成绩和表现做一份家长会材料。” & \texttt{class-artifacts} + 查询 Skills & 只取必要字段生成 Office 草稿，登记来源记录后再交付 \\
    \bottomrule
  \end{tabular}
\end{table}

在最后一个例子中，文件生成不是数据写入的替代品。产物需要带有来源实体和生成时间；如果用户修改了主数据，应重新生成或明确标记旧产物仍对应哪个版本。

\clearpage
\section{工程校验与证据边界}

仓库把协议和安装结构做成可以重复运行的检查，而不是只在 README 中描述约定：

\begin{itemize}
  \item \texttt{scripts/validate\_schema.py} 校验数据协议实例，\texttt{scripts/contract\_test.py} 检查操作和安全约束；
  \item Codex 插件清单位于 \texttt{.codex-plugin/plugin.json}，每个 \texttt{skills/*/SKILL.md} 都可以单独做格式检查；
  \item \texttt{tools/setup\_doctor.py} 报告本地工具和连接器状态，不能替代真实账号授权；
  \item \texttt{references/capability-matrix.md} 明确区分飞书的优先适配与 Notion、Obsidian 的映射和 fixture 状态。
\end{itemize}

这些检查只能证明仓库结构、协议和 fixture 的一致性，不能证明某个学校已经完成部署，也不能推出学生成绩、班级氛围或教师负担发生了统计意义上的变化。涉及真实学生资料时，仍需要学校的授权、最小化采集、访问控制和人工复核。

\section{局限与后续工作}

当前版本的主要边界如下：

\begin{enumerate}
  \item 飞书、Notion 和 Obsidian 的连接器权限、API 版本和可用 CLI 由运行环境决定，仓库不代替这些配置；
  \item 权限角色、字段级审计、家长可见视图和冲突解决仍需结合实际学校流程设计；
  \item 排座、值日和家长会材料已有规则入口，但复杂模板、版式主题和多轮人工修改尚未形成完整的模板管理层；
  \item 当前协议支持显式导入、追加、更新、查询和导出，跨后端自动双向同步、自动预警和效果评估不在默认范围内。
\end{enumerate}

后续可以围绕三个方向推进：为 Notion 和 Obsidian 增加最小可用初始化及更多 live connector 测试；补充角色与审计模型；在不把 Office 文件当作主数据库的前提下，完善可复用的产物模板和版本比较。

\section{结论}

班主任.Skill v3 将班级治理任务组织成可安装的多 Skill 插件：业务 Skills 负责记录、查询、成长跟进、班务安排和文件生成，适配 Skills 负责把统一协议投影到不同后端。稳定 ID、版本、来源、敏感级别以及预览确认机制，为日常使用提供了可检查的边界。当前版本已经把协议、插件结构和三类后端映射放入同一仓库；实时能力和教育效果仍应依据具体连接器、授权和后续实证逐项验证。

\end{document}
